%\SweaveUTF8



\setkeys{Gin}{width = \textwidth}

\begin{Problem}[Stratified simple random sampling]
    We again use the data set \code{env.dat}. Assume that membership in
    industry class 31 was known from the sampling frame and that the population
    was stratified by this variable before sampling. Within each stratum, a
    simple random sample without replacement was drawn.
    The population consists of two strata:

    \vspace{0.5em}
    \begin{center}
        \begin{tabular}{ll}
            \toprule
            \addlinespace
            Stratum & Population size\\
            \addlinespace
            \midrule
            \addlinespace
            $U_{31}$ (\code{code == 1}) & $N_{31} = 460$\\
            \addlinespace
            $U_{\mathrm{rest}}$ (\code{code == 0}) &
                $N_{\mathrm{rest}} = 3{,}310$\\
            \bottomrule
        \end{tabular}
    \end{center}

    \vspace{0.5em}

    \begin{Rhint}
        The stratum sample sizes are available in \code{sepe.dat}. As explained
        in Problem Set 2, create a survey design object for stratified sampling
        and name it \code{dn_sepe}.
    \end{Rhint}

\begin{Subtask}
    %---------------------------------------------------------------------------
    \item Estimate the population total of environmental-protection
        expenditures using \code{svytotal()}. Also report a 90\,\% confidence
        interval using

\begin{Scode}
confint(object, parm, level = 0.95, ...)
\end{Scode}

    \begin{Solution}
        Load the data and define the survey design:


\begin{Schunk}
\begin{Sinput}
> library("survey")
> sepe <- read.table("https://stat.ethz.ch/Teaching/Datasets/cas-das/sepe.dat",
                     header = TRUE)
\end{Sinput}
\end{Schunk}

\begin{Schunk}
\begin{Sinput}
> sepe$Nph <- ifelse(sepe$code == 1, 460, 3310)
> dn_sepe <- svydesign(ids = ~1, strata = ~code, fpc = ~Nph, data = sepe)
\end{Sinput}
\end{Schunk}

        The estimated total and its 90\,\% confidence interval are

\begin{Schunk}
\begin{Sinput}
> t_y <- svytotal(~expend, dn_sepe)
> t_y
\end{Sinput}
\begin{Soutput}
           total      SE
expend 317824493 7719366
\end{Soutput}
\begin{Sinput}
> confint(t_y, level = 0.90)
\end{Sinput}
\begin{Soutput}
             5 %      95 %
expend 305127266 330521720
\end{Soutput}
\end{Schunk}
    \end{Solution}

    %---------------------------------------------------------------------------
    \item Estimate the mean environmental-protection expenditures in each
        stratum using \code{svyby()}.

        \begin{Rhint}
            The function \code{svyby(formula, by, design, FUN, ...)} applies
            the survey statistic \code{FUN} to the variables specified by
            \code{formula}, separately for the groups defined by \code{by},
            using the survey design \code{design}. Additional arguments
            can be passed through \code{...}.
        \end{Rhint}

        Next, test whether the two stratum means are equal. For comparison,
        use both

\begin{Scode}
t.test(expend ~ code, data = sepe, var.equal = TRUE)
\end{Scode}

        and the design-based test

\begin{Scode}
svyttest(expend ~ code, dn_sepe)
\end{Scode}

        The first is the classical pooled two-sample $t$-test and ignores the
        sampling design. The second uses the variance implied by the survey
        design, including the finite population corrections.

    \begin{Solution}
        The estimated stratum means are

\begin{Schunk}
\begin{Sinput}
> svyby(~expend, ~code, dn_sepe, svymean)
\end{Sinput}
\begin{Soutput}
  code    expend       se
0    0  80308.39 2173.576
1    1 113051.55 6082.014
\end{Soutput}
\end{Schunk}

        The classical pooled two-sample $t$-test is

\begin{Schunk}
\begin{Sinput}
> t.test(expend ~ code, data = sepe, var.equal = TRUE)
\end{Sinput}
\begin{Soutput}
	Two Sample t-test

data:  expend by code
t = -5.2217, df = 798, p-value = 2.264e-07
alternative hypothesis: true difference in means between group 0 and group 1 is not equal to 0
95 percent confidence interval:
 -45052.07 -20434.25
sample estimates:
mean in group 0 mean in group 1 
       80308.39       113051.55 
\end{Soutput}
\end{Schunk}

        The corresponding design-based test is

\begin{Schunk}
\begin{Sinput}
> svyttest(expend ~ code, dn_sepe)
\end{Sinput}
\begin{Soutput}
	Design-based t-test

data:  expend ~ code
t = 5.0696, df = 797, p-value = 4.96e-07
alternative hypothesis: true difference in mean is not equal to 0
95 percent confidence interval:
 20065.01 45421.31
sample estimates:
difference in mean 
          32743.16 
\end{Soutput}
\end{Schunk}

        The two tests need not give the same standard error or $p$-value.
        \code{t.test()} treats the observations as independent samples from
        two populations and pools their variances, whereas \code{svyttest()}
        uses a design-based variance estimate from \code{dn_sepe}.
    \end{Solution}

    %---------------------------------------------------------------------------
    \item Let $y_i$ denote environmental-protection expenditures
        (\code{expend}) and $x_i$ the number of full-time employees
        (\code{empl}). Two different summaries are commonly considered:
        \begin{equation*}
            r_{AR}
            = \frac{1}{N_+}\sum_{i \in U_+}\frac{y_i}{x_i},
            \qquad
            r_{RA}
            = \frac{\sum_{i \in U} y_i}{\sum_{i \in U} x_i},
        \end{equation*}
        where $U_+ = \{i \in U:x_i>0\}$ and $N_+ = |U_+|$. The first is the
        average of ratios; the second is the ratio of totals
        (equivalently, the ratio of averages).

        Estimate $r_{RA}$ using \code{svyratio()} and the design object
        \code{dn_sepe}. Then repeat the calculation with
        \code{separate = TRUE}. What changes?

        \begin{Rhint}
            Use

\begin{Scode}
svyratio(numerator = formula, denominator, design, ...)
\end{Scode}

            The default estimates one combined ratio. With \code{separate =
            TRUE}, ratios are estimated separately within the strata.
        \end{Rhint}

    \begin{Solution}
        The combined ratio estimator is
        \begin{equation*}
            \widehat r_{RA}
            = \frac{\widehat t_y}{\widehat t_x}.
        \end{equation*}

\begin{Schunk}
\begin{Sinput}
> r_RA <- svyratio(~expend, ~empl, dn_sepe)
> r_RA
\end{Sinput}
\begin{Soutput}
Ratio estimator: svyratio.survey.design2(~expend, ~empl, dn_sepe)
Ratios=
           empl
expend 450.1731
SEs=
           empl
expend 10.85506
\end{Soutput}
\end{Schunk}

        Thus, the estimated ratio is
        450.2, with estimated standard error
        10.9.
        With \code{separate = TRUE}, \code{svyratio()} instead reports one
        ratio for each stratum:

\begin{Schunk}
\begin{Sinput}
> r_RA_sep <- svyratio(~expend, ~empl, dn_sepe, separate = TRUE)
> r_RA_sep
\end{Sinput}
\begin{Soutput}
Stratified ratio estimate: svyratio.survey.design2(~expend, ~empl, dn_sepe, separate = TRUE)
Ratio estimator: Stratum == 0L
Ratios=
           empl
expend 461.3302
SEs=
           empl
expend 12.77544
Ratio estimator: Stratum == 1L
Ratios=
           empl
expend 400.6449
SEs=
           empl
expend 18.37954
\end{Soutput}
\end{Schunk}

        Hence, \code{separate = TRUE} does not produce a second single
        estimate of the overall ratio; it produces the within-stratum ratios
        $\widehat t_{yh}/\widehat t_{xh}$.
    \end{Solution}

    %---------------------------------------------------------------------------
    \item Convert \code{dn_sepe} to a replicate-weight design for the
        Rao--Wu $(n-1)$ bootstrap with $B = 200$ replicates:

\begin{Scode}
set.seed(11)
dn_sepe_boot <- as.svrepdesign(dn_sepe, type = "subbootstrap",
                               replicates = 200)
\end{Scode}

        Re-estimate the combined ratio $r_{RA}$ from part c) using this
        replicate design. Also report its 95\,\% confidence interval.

    \begin{Solution}
        Create the replicate-weight design:

\begin{Schunk}
\begin{Sinput}
> set.seed(11)
> dn_sepe_boot <- as.svrepdesign(dn_sepe, type = "subbootstrap",
                                replicates = 200)
\end{Sinput}
\end{Schunk}

        Estimate the ratio:

\begin{Schunk}
\begin{Sinput}
> r_RA_boot <- svyratio(~expend, ~empl, dn_sepe_boot)
> r_RA_boot
\end{Sinput}
\begin{Soutput}
Ratio estimator: svyratio.svyrep.design(~expend, ~empl, dn_sepe_boot)
Ratios=
           empl
expend 450.1731
SEs=
         [,1]
[1,] 12.14926
\end{Soutput}
\end{Schunk}

        The point estimate is the same ratio estimator as before, while its
        standard error is now estimated from the bootstrap replicate weights.
        The 95\,\% confidence interval is

\begin{Schunk}
\begin{Sinput}
> confint(r_RA_boot)
\end{Sinput}
\begin{Soutput}
              2.5 %   97.5 %
expend/empl 426.361 473.9853
\end{Soutput}
\end{Schunk}
    \end{Solution}
\end{Subtask}
\end{Problem}
